\clearpage
\begin{trapbox}
\textbf{误区一：按字母大小套公式，交换两个分母。}
\begin{center}\begin{tikzpicture}[font=\small]
\draw[black!30,rounded corners=2pt] (0,0) rectangle (14.4,8.2);
\node[anchor=west,font=\small\bfseries] at (.5,7.72) {认分母，不认字母大小：横长椭圆与纵长椭圆是不同对象};
% MODEL 9 4 0 1 0
\begin{scope}[shift={(4,3.6)},scale=0.8]
\draw[blue,very thick] (0,0) ellipse (3.0000000000 and 2.0000000000);
\draw[-Stealth,black!45] (-3.4,0)--(3.4,0) node[right] {$x$};\draw[-Stealth,black!45] (0,-3.1)--(0,3.1) node[above] {$y$};

\end{scope}
\begin{scope}[shift={(10.4,3.6)},scale=0.8]
\draw[blue,very thick] (0,0) ellipse (2.0000000000 and 3.0000000000);
\draw[-Stealth,black!45] (-3.4,0)--(3.4,0) node[right] {$x$};\draw[-Stealth,black!45] (0,-3.1)--(0,3.1) node[above] {$y$};

\end{scope}
\node at (4,6.9) {$x^2/9+y^2/4=1$};\node at (10.4,6.9) {$x^2/4+y^2/9=1$};\node[anchor=west] at (.8,.6) {左右单位相同：交换分母会改变曲线，不能随意对换};
\end{tikzpicture}\end{center}

左图 $\alpha=9,\beta=4$，右图 $\alpha=4,\beta=9$，曲线方向已改变。
\textbf{正确替代：} 先看 $x^2,y^2$ 下方各是什么数，再代 $\alpha,\beta$；
它们是平方值，不是半轴长。无需另背两套焦点轴公式。

\textbf{误区二：未检查消元条件，就用判别式判位置。}
竖直弦图已说明：$B=0$ 时，横坐标相同并不说明相切。
\textbf{正确替代：} 消 $y$ 检查 $B\ne0$，消 $x$ 检查 $A\ne0$；
或者用对所有方向都有效的 $H=D-C^2$。要有两个不同交点，必须严格满足 $H>0$。

\tcbbreak
\textbf{误区三：斜截式移项时符号写错。}
\begin{center}\begin{tikzpicture}[font=\small]
\draw[black!30,rounded corners=2pt] (0,0) rectangle (14.4,8.2);
\node[anchor=west,font=\small\bfseries] at (.5,7.72) {移项错一个符号，就换了几何对象：先核对直线一般式};
% MODEL 8 4 1 1 -1
\begin{scope}[shift={(7.2,3.65)},scale=1.08]
\draw[blue,very thick] (0,0) ellipse (2.8284271247 and 2.0000000000);
\draw[-Stealth,black!45] (-3.4,0)--(3.4,0) node[right] {$x$};\draw[-Stealth,black!45] (0,-2.3)--(0,2.3) node[above] {$y$};
\draw[orange!85!black,thick] (-1.2,2.2)--(2.6,-1.6);\draw[red!75!black,thick,dashed] (-1.2,-2.2)--(2.6,1.6);
\end{scope}
\node[anchor=west,text=orange!85!black] at (.8,6.8) {原式 $y=-x+1$：$x+y-1=0$};\node[anchor=west,text=red!75!black] at (.8,.6) {错写 $-x+y+1=0$，得到的是另一条直线};
\end{tikzpicture}\end{center}

一般地，$y=kx+m$ 可写成 $kx-y+m=0$，也可整体乘 $-1$；
但不能只改其中某一项的符号。图中错式已不是原直线。
\textbf{正确替代：} 代一个已知点回原式检查；若整式同时乘非零数，表示的才是同一条线。
在 $A,B,C$ 同乘 $s\ne0$ 时，$D,H$ 同乘 $s^2$，中点与弦长不变。

\textbf{误区四：混用约分前后的判别式，或把根差当弦长。}
若二次方程整体乘 $s\ne0$，判别式乘 $s^2$，二次项系数乘 $s$，
因此 $\sqrt\Delta/|p|$ 不变。只改分母或只改判别式都会算错。
\textbf{正确替代：} 系数与判别式必须来自同一个方程；横坐标根差还须乘 $\sqrt{1+k^2}$。
竖直线没有有限斜率，直接用纵坐标差或通用弦长式。

\tcbbreak
\textbf{误区五：把椭圆当作圆，套用垂径定理。}
\begin{center}\begin{tikzpicture}[font=\small]
\draw[black!30,rounded corners=2pt] (0,0) rectangle (14.4,8.2);
\node[anchor=west,font=\small\bfseries] at (.5,7.72) {椭圆的弦中点一般不是垂足：$OM$ 不必垂直于弦};
% MODEL 9 1 1 2 -2
\begin{scope}[shift={(7,3.6)},scale=1.4]
\draw[blue,very thick] (0,0) ellipse (3.0000000000 and 1.0000000000);
\coordinate (P) at (0.0000000000,1.0000000000);
\coordinate (Q) at (2.7692307692,-0.3846153846);
\coordinate (M) at (1.3846153846,0.3076923077);
\coordinate (O) at (0.0000000000,0.0000000000);
\coordinate (N) at (0.4000000000,0.8000000000);
\draw[orange!85!black,very thick] (P)--(Q);\draw[purple,thick] (O)--(M);\draw[green!45!black,thick,dashed] (O)--(N);\pic[draw,angle radius=3mm] {right angle=O--N--Q};\fill[black] (O) circle (1.8pt) node[below left=6pt] {$O$};
\fill[black] (P) circle (1.8pt) node[above left=6pt] {$P$};
\fill[black] (Q) circle (1.8pt) node[below right=6pt] {$Q$};
\fill[black] (M) circle (1.8pt) node[below=13pt] {$M$};
\fill[black] (N) circle (1.8pt) node[below left=7pt] {$N$};

\end{scope}
\node[anchor=west] at (.8,6.7) {$x^2/9+y^2=1,\quad x+2y-2=0$};\node[anchor=west] at (.8,1.0) {弦中点 $M=(18/13,8/13)$};\node[anchor=west] at (.8,.5) {垂足 $N=(2/5,4/5)$；$M\ne N$，不能套圆的垂径定理};
\end{tikzpicture}\end{center}

本图 $D=13,H=9>0$，公式给出弦中点 $M=(18/13,8/13)$；
从原点到直线的垂足却为 $N=(2/5,4/5)$。
$M\ne N$，所以错误的“作垂线就找中点”会找到另一个点。
\textbf{正确替代：} 椭圆用中点公式，或由两点平方差得到
$Bx_M/\alpha-Ay_M/\beta=0$，而不是擅自令 $OM\perp PQ$。
只有特殊弦或圆的情形才能出现垂直重合，不能当作一般性质。
\end{trapbox}
